\documentclass[aspectratio=169,11pt]{beamer}

% ============================================
% PACKAGES
% ============================================
\usepackage[utf8]{inputenc}
\usepackage[T1]{fontenc}
\usepackage[ngerman]{babel}
\usepackage{lmodern}
\usepackage{tcolorbox}
\tcbuselibrary{skins, hooks}
\usepackage{listings}
\usepackage{xcolor}
\usepackage{fontawesome5}
\usepackage{tikz}
\usetikzlibrary{shapes, arrows.meta}

% Modern Font
\IfFileExists{FiraSans.sty}{\usepackage[sfdefault]{FiraSans}}{}
\renewcommand{\familydefault}{\sfdefault}

% ============================================
% COLORS & DESIGN SYSTEM
% ============================================
\definecolor{BgColor}{HTML}{FAFAFA}
\definecolor{Primary}{HTML}{2B3A42}
\definecolor{Accent}{HTML}{E84A5F}
\definecolor{Secondary}{HTML}{3F5765}
\definecolor{CodeBg}{HTML}{F0F0F0}
\definecolor{Success}{HTML}{28A745}

\setbeamercolor{background canvas}{bg=BgColor}
\setbeamercolor{title}{fg=Primary}
\setbeamercolor{frametitle}{bg=Primary, fg=white}
\setbeamercolor{structure}{fg=Accent}
\setbeamercolor{normal text}{fg=Primary}
\setbeamercolor{itemize item}{fg=Accent}

\setbeamertemplate{navigation symbols}{}
\setbeamertemplate{footline}{%
  \leavevmode%
  \hbox{%
  \begin{beamercolorbox}[wd=.333333\paperwidth,ht=2.25ex,dp=1ex,center]{author in head/foot}%
    \usebeamerfont{author in head/foot}\tiny Falko Himmel
  \end{beamercolorbox}%
  \begin{beamercolorbox}[wd=.333333\paperwidth,ht=2.25ex,dp=1ex,center]{title in head/foot}%
    \usebeamerfont{title in head/foot}\tiny Grundlagen der Praktischen Informatik
  \end{beamercolorbox}%
  \begin{beamercolorbox}[wd=.333333\paperwidth,ht=2.25ex,dp=1ex,right]{date in head/foot}%
    \usebeamerfont{date in head/foot}\insertframenumber{} / \inserttotalframenumber\hspace*{2ex} 
  \end{beamercolorbox}}%
  \vskip0pt%
}

% ============================================
% CUSTOM BOXES
% ============================================
\tcbset{
    modernbox/.style={
        enhanced,
        boxrule=0pt,
        frame hidden,
        colback=white,
        coltitle=white,
        colbacktitle=Secondary,
        fonttitle=\bfseries,
        drop lifted shadow,
        arc=4pt,
        toptitle=2mm,
        bottomtitle=2mm,
        left=2mm,
        right=2mm
    },
    alertbox/.style={
        modernbox,
        colbacktitle=Accent
    }
}

% ============================================
% CODE LISTINGS
% ============================================
\lstdefinelanguage{Haskell}{
  keywords={class, data, deriving, do, else, if, in, import, let, module, of, then, type, where, case, error, instance},
  keywordstyle=\color{Accent}\bfseries,
  identifierstyle=\color{Primary},
  sensitive=true,
  comment=[l]{--},
  commentstyle=\color{gray}\ttfamily\itshape,
  stringstyle=\color{teal}\ttfamily,
  morestring=[b]"
}

\lstset{
    language=Haskell,
    backgroundcolor=\color{CodeBg},
    basicstyle=\ttfamily\scriptsize,
    breaklines=true,
    frame=single,
    rulecolor=\color{CodeBg},
    frameround=tttt,
    numbers=left,
    numberstyle=\tiny\color{gray},
    xleftmargin=15pt,
    tabsize=4,
    showstringspaces=false
}

% ============================================
% TITLE PAGE
% ============================================
\title{\textbf{Tutorium 8}}
\subtitle{Abstrakte Datenstrukturen \& Hash-Maps}
\institute{Grundlagen der Praktischen Informatik}
\author{}
\date{\today}

\begin{document}

\begin{frame}[plain]
    \titlepage
\end{frame}

\begin{frame}{Hash-Maps: Visualisierung}
    \begin{center}
    \begin{tikzpicture}
        \node (key) at (0,2) {\textbf{Key: ,,Apfel``}};
        \node[draw, rounded corners, fill=Secondary!20] (hash) at (3,2) {$h(k) = f(k) \pmod m$};
        \node (index) at (6,2) {\textbf{Index: 2}};
        
        \draw[->, thick] (key) -- (hash);
        \draw[->, thick] (hash) -- (index);
        
        \draw[thick] (5,-1) grid (7,1.5);
        \node at (6, 1.25) {0};
        \node at (6, 0.75) {1};
        \node[fill=Accent!30, minimum width=1.9cm, minimum height=0.4cm] at (6, 0.25) {Value};
        \node at (4.5, 0.25) {\footnotesize $\leftarrow$ Index 2};
        \node at (6, -0.25) {3};
        \node at (6, -0.75) {4};
        
        \draw[->, thick, Accent] (index) -- (6, 0.5);
    \end{tikzpicture}
    \end{center}
\end{frame}


\begin{frame}{Inhaltsverzeichnis}
    \tableofcontents
\end{frame}

\begin{frame}{Hash-Maps: Visualisierung}
    \begin{center}
    \begin{tikzpicture}
        \node (key) at (0,2) {\textbf{Key: ,,Apfel``}};
        \node[draw, rounded corners, fill=Secondary!20] (hash) at (3,2) {$h(k) = f(k) \pmod m$};
        \node (index) at (6,2) {\textbf{Index: 2}};
        
        \draw[->, thick] (key) -- (hash);
        \draw[->, thick] (hash) -- (index);
        
        \draw[thick] (5,-1) grid (7,1.5);
        \node at (6, 1.25) {0};
        \node at (6, 0.75) {1};
        \node[fill=Accent!30, minimum width=1.9cm, minimum height=0.4cm] at (6, 0.25) {Value};
        \node at (4.5, 0.25) {\footnotesize $\leftarrow$ Index 2};
        \node at (6, -0.25) {3};
        \node at (6, -0.75) {4};
        
        \draw[->, thick, Accent] (index) -- (6, 0.5);
    \end{tikzpicture}
    \end{center}
\end{frame}


% ============================================
% STACKS & QUEUES
% ============================================
\section{Stacks und Queues}

\begin{frame}{Abstrakte Datenstrukturen (ADTs)}
    Stacks und Queues definieren, in welcher Reihenfolge Elemente eingefügt und entnommen werden.
    
    \begin{columns}[T]
        \begin{column}{0.48\textwidth}
            \begin{tcolorbox}[modernbox, title=\faLayerGroup\ Stack (LIFO)]
                \textbf{L}ast \textbf{I}n, \textbf{F}irst \textbf{O}ut
                \begin{itemize}
                    \item \texttt{push}: Element oben drauflegen
                    \item \texttt{pop}: Oberstes Element entfernen
                \end{itemize}
                \textit{Beispiel: Undo-Historie, Tiefensuche.}
            \end{tcolorbox}
        \end{column}
        \begin{column}{0.48\textwidth}
            \begin{tcolorbox}[modernbox, title=\faUsers\ Queue (FIFO)]
                \textbf{F}irst \textbf{I}n, \textbf{F}irst \textbf{O}ut
                \begin{itemize}
                    \item \texttt{enqueue}: Element hinten anstellen
                    \item \texttt{dequeue}: Vorderstes Element entfernen
                \end{itemize}
                \textit{Beispiel: Warteschlangen, Breitensuche.}
            \end{tcolorbox}
        \end{column}
    \end{columns}
\end{frame}

\begin{frame}{Hash-Maps: Visualisierung}
    \begin{center}
    \begin{tikzpicture}
        \node (key) at (0,2) {\textbf{Key: ,,Apfel``}};
        \node[draw, rounded corners, fill=Secondary!20] (hash) at (3,2) {$h(k) = f(k) \pmod m$};
        \node (index) at (6,2) {\textbf{Index: 2}};
        
        \draw[->, thick] (key) -- (hash);
        \draw[->, thick] (hash) -- (index);
        
        \draw[thick] (5,-1) grid (7,1.5);
        \node at (6, 1.25) {0};
        \node at (6, 0.75) {1};
        \node[fill=Accent!30, minimum width=1.9cm, minimum height=0.4cm] at (6, 0.25) {Value};
        \node at (4.5, 0.25) {\footnotesize $\leftarrow$ Index 2};
        \node at (6, -0.25) {3};
        \node at (6, -0.75) {4};
        
        \draw[->, thick, Accent] (index) -- (6, 0.5);
    \end{tikzpicture}
    \end{center}
\end{frame}



\begin{frame}{Visualisierung}
    \begin{columns}[T]
        \begin{column}{0.48\textwidth}
            \begin{center}
            \textbf{Stack (LIFO)}\\
            \vspace{0.3cm}
            \begin{tikzpicture}
                \draw[thick] (0,0) -- (0,2);
                \draw[thick] (1.5,0) -- (1.5,2);
                \draw[thick] (0,0) -- (1.5,0);
                \draw[fill=Accent!20] (0.1,0.1) rectangle (1.4,0.6) node[pos=.5] {1};
                \draw[fill=Accent!40] (0.1,0.7) rectangle (1.4,1.2) node[pos=.5] {2};
                \draw[->, thick, Accent] (2, 2) to[bend left] node[right] {\tiny push 3} (0.75, 1.4);
                \draw[->, thick, Secondary] (0.75, 1.4) to[bend left] node[left] {\tiny pop} (-0.5, 2);
            \end{tikzpicture}
            \end{center}
        \end{column}
        \begin{column}{0.48\textwidth}
            \begin{center}
            \textbf{Queue (FIFO)}\\
            \vspace{0.3cm}
            \begin{tikzpicture}
                \draw[thick] (0,0) -- (3,0);
                \draw[thick] (0,1) -- (3,1);
                \draw[fill=Success!20] (0.2,0.1) rectangle (0.8,0.9) node[pos=.5] {A};
                \draw[fill=Success!40] (1.0,0.1) rectangle (1.6,0.9) node[pos=.5] {B};
                \draw[fill=Success!60] (1.8,0.1) rectangle (2.4,0.9) node[pos=.5] {C};
                \draw[<-, thick, Success] (2.6, 0.5) -- (3.5, 0.5) node[right] {\tiny enqueue};
                \draw[->, thick, Secondary] (-0.8, 0.5) -- (0, 0.5) node[left] {\tiny dequeue};
            \end{tikzpicture}
            \end{center}
        \end{column}
    \end{columns}
\end{frame}

\begin{frame}{Hash-Maps: Visualisierung}
    \begin{center}
    \begin{tikzpicture}
        \node (key) at (0,2) {\textbf{Key: ,,Apfel``}};
        \node[draw, rounded corners, fill=Secondary!20] (hash) at (3,2) {$h(k) = f(k) \pmod m$};
        \node (index) at (6,2) {\textbf{Index: 2}};
        
        \draw[->, thick] (key) -- (hash);
        \draw[->, thick] (hash) -- (index);
        
        \draw[thick] (5,-1) grid (7,1.5);
        \node at (6, 1.25) {0};
        \node at (6, 0.75) {1};
        \node[fill=Accent!30, minimum width=1.9cm, minimum height=0.4cm] at (6, 0.25) {Value};
        \node at (4.5, 0.25) {\footnotesize $\leftarrow$ Index 2};
        \node at (6, -0.25) {3};
        \node at (6, -0.75) {4};
        
        \draw[->, thick, Accent] (index) -- (6, 0.5);
    \end{tikzpicture}
    \end{center}
\end{frame}


\begin{frame}[fragile]{Umsetzung in Haskell (Typklassen)}
    Wir können das Verhalten über Typklassen definieren.
    
    \begin{tcolorbox}[modernbox, title=Stack-Implementierung mit Listen]
\begin{lstlisting}
class Stack s where
    push   :: a -> s a -> s a 
    pop    :: s a -> (a, s a)
    sEmpty :: s a -> Bool

instance Stack [] where
    push       = (:)
    pop []     = error "Stack is empty!"
    pop (x:xs) = (x, xs)
    sEmpty     = null
\end{lstlisting}
    \end{tcolorbox}
\end{frame}

\begin{frame}{Hash-Maps: Visualisierung}
    \begin{center}
    \begin{tikzpicture}
        \node (key) at (0,2) {\textbf{Key: ,,Apfel``}};
        \node[draw, rounded corners, fill=Secondary!20] (hash) at (3,2) {$h(k) = f(k) \pmod m$};
        \node (index) at (6,2) {\textbf{Index: 2}};
        
        \draw[->, thick] (key) -- (hash);
        \draw[->, thick] (hash) -- (index);
        
        \draw[thick] (5,-1) grid (7,1.5);
        \node at (6, 1.25) {0};
        \node at (6, 0.75) {1};
        \node[fill=Accent!30, minimum width=1.9cm, minimum height=0.4cm] at (6, 0.25) {Value};
        \node at (4.5, 0.25) {\footnotesize $\leftarrow$ Index 2};
        \node at (6, -0.25) {3};
        \node at (6, -0.75) {4};
        
        \draw[->, thick, Accent] (index) -- (6, 0.5);
    \end{tikzpicture}
    \end{center}
\end{frame}


% ============================================
% HASH-MAPS
% ============================================
\section{Hash-Maps}

\begin{frame}{Hash-Maps: Grundlagen}
    Eine \textbf{Hash-Map} ordnet Schlüssel (Keys) auf Werte (Values) zu. Die Position in der Tabelle wird durch Hashing berechnet. 
    
    \begin{tcolorbox}[modernbox, title=Hash-Funktion \& Lastfaktor]
        \begin{itemize}
            \item \textbf{Kongruenzmethode:} $h(k) = f(k) \pmod m$\\
                  $m$ wird idealerweise als Primzahl gewählt.
            \item \textbf{Lastfaktor ($\alpha$):} Verhältnis zwischen belegten und verfügbaren Plätzen ($\alpha = \frac{n}{m}$).
        \end{itemize}
    \end{tcolorbox}
    
    \textbf{Laufzeit:} Im Mittel $\mathcal{O}(1)$ für Suche, Einfügen und Löschen!
\end{frame}

\begin{frame}{Hash-Maps: Visualisierung}
    \begin{center}
    \begin{tikzpicture}
        \node (key) at (0,2) {\textbf{Key: ,,Apfel``}};
        \node[draw, rounded corners, fill=Secondary!20] (hash) at (3,2) {$h(k) = f(k) \pmod m$};
        \node (index) at (6,2) {\textbf{Index: 2}};
        
        \draw[->, thick] (key) -- (hash);
        \draw[->, thick] (hash) -- (index);
        
        \draw[thick] (5,-1) grid (7,1.5);
        \node at (6, 1.25) {0};
        \node at (6, 0.75) {1};
        \node[fill=Accent!30, minimum width=1.9cm, minimum height=0.4cm] at (6, 0.25) {Value};
        \node at (4.5, 0.25) {\footnotesize $\leftarrow$ Index 2};
        \node at (6, -0.25) {3};
        \node at (6, -0.75) {4};
        
        \draw[->, thick, Accent] (index) -- (6, 0.5);
    \end{tikzpicture}
    \end{center}
\end{frame}



\begin{frame}{Hash-Maps: Visualisierung}
    \begin{center}
    \begin{tikzpicture}
        \node (key) at (0,2) {\textbf{Key: ,,Apfel``}};
        \node[draw, rounded corners, fill=Secondary!20] (hash) at (3,2) {$h(k) = f(k) \pmod m$};
        \node (index) at (6,2) {\textbf{Index: 2}};
        
        \draw[->, thick] (key) -- (hash);
        \draw[->, thick] (hash) -- (index);
        
        \draw[thick] (5,-1) grid (7,1.5);
        \node at (6, 1.25) {0};
        \node at (6, 0.75) {1};
        \node[fill=Accent!30, minimum width=1.9cm, minimum height=0.4cm] at (6, 0.25) {Value};
        \node at (4.5, 0.25) {\footnotesize $\leftarrow$ Index 2};
        \node at (6, -0.25) {3};
        \node at (6, -0.75) {4};
        
        \draw[->, thick, Accent] (index) -- (6, 0.5);
    \end{tikzpicture}
    \end{center}
\end{frame}

\begin{frame}{Kollisionen}
    \begin{tcolorbox}[alertbox, title=Kollision]
        Eine Kollision liegt vor, wenn $h(k) = h(k')$ für zwei verschiedene Schlüssel $k \neq k'$.
    \end{tcolorbox}
    
    \textbf{Strategien zur Behandlung:}
    \begin{itemize}
        \item \textbf{Separate Chaining:} Eine Liste / Bucket pro Tabellenfeld.
        \item \textbf{Offene Adressierung:} Bei Belegung weicht man auf andere Indizes aus.
        \begin{itemize}
            \item Lineares Sondieren ($+1, +2, +3 \dots$)
            \item Quadratisches Sondieren ($+1^2, +2^2, +3^2 \dots$)
            \item Double Hashing (Zweite Hash-Funktion berechnet Sprungweite)
        \end{itemize}
    \end{itemize}
\end{frame}

\begin{frame}{Hash-Maps: Visualisierung}
    \begin{center}
    \begin{tikzpicture}
        \node (key) at (0,2) {\textbf{Key: ,,Apfel``}};
        \node[draw, rounded corners, fill=Secondary!20] (hash) at (3,2) {$h(k) = f(k) \pmod m$};
        \node (index) at (6,2) {\textbf{Index: 2}};
        
        \draw[->, thick] (key) -- (hash);
        \draw[->, thick] (hash) -- (index);
        
        \draw[thick] (5,-1) grid (7,1.5);
        \node at (6, 1.25) {0};
        \node at (6, 0.75) {1};
        \node[fill=Accent!30, minimum width=1.9cm, minimum height=0.4cm] at (6, 0.25) {Value};
        \node at (4.5, 0.25) {\footnotesize $\leftarrow$ Index 2};
        \node at (6, -0.25) {3};
        \node at (6, -0.75) {4};
        
        \draw[->, thick, Accent] (index) -- (6, 0.5);
    \end{tikzpicture}
    \end{center}
\end{frame}


% ============================================
% LIVE DEMO
% ============================================
\section{Praxis}

\begin{frame}[plain]
    \begin{center}
        \Huge \textbf{Live Demo} \\
        \vspace{0.5cm}
        \large \faLaptopCode~ \texttt{02\_Abstrakte\_Datenstrukturen.hs} \\
        \vspace{0.3cm}
        \normalsize Gemeinsame Besprechung von Hash-Maps \& Klausuraufgabe (2023).
\vspace{0.5cm}\begin{center}\href{https://moodle.uni-ulm.de/pluginfile.php/1348606/mod_folder/content/0/2023_1_exam.pdf}{\beamerbutton{Altklausur 2023\_1 auf Moodle öffnen}}\end{center}
    \end{center}
\end{frame}

\begin{frame}{Hash-Maps: Visualisierung}
    \begin{center}
    \begin{tikzpicture}
        \node (key) at (0,2) {\textbf{Key: ,,Apfel``}};
        \node[draw, rounded corners, fill=Secondary!20] (hash) at (3,2) {$h(k) = f(k) \pmod m$};
        \node (index) at (6,2) {\textbf{Index: 2}};
        
        \draw[->, thick] (key) -- (hash);
        \draw[->, thick] (hash) -- (index);
        
        \draw[thick] (5,-1) grid (7,1.5);
        \node at (6, 1.25) {0};
        \node at (6, 0.75) {1};
        \node[fill=Accent!30, minimum width=1.9cm, minimum height=0.4cm] at (6, 0.25) {Value};
        \node at (4.5, 0.25) {\footnotesize $\leftarrow$ Index 2};
        \node at (6, -0.25) {3};
        \node at (6, -0.75) {4};
        
        \draw[->, thick, Accent] (index) -- (6, 0.5);
    \end{tikzpicture}
    \end{center}
\end{frame}


\end{document}
